\documentclass[12pt, letterpaper]{article} 

\usepackage{amsmath, amssymb, amsthm, enumerate}

\title{MATH 189 HW 1}
\author{Dani Nguyen}
\date{Winter 2025}

\begin{document}
\maketitle

\noindent \textbf{Question 1}
\begin{enumerate}[a.]
    \item Regression, $n = 500$, $p = 3$
    \item Categorical, $n =20$, $p = 14$
    \item Regression, $n = 52$, $p = 3$
\end{enumerate}

\noindent \textbf{Question 2}
\begin{enumerate}[a.]
    \item Obs. 1: 2 \\
          Obs. 2: 2 \\
          Obs. 3: $\sqrt{1^2+3^2} = \sqrt{10} \approx 3.16$ \\
          Obs. 4: $\sqrt{1^2+2^2} = \sqrt{5} \approx 2.23$ \\
          Obs. 5: $\sqrt{(-1)^2 + 1^2} = \sqrt{2} \approx 1.41$ \\
          Obs. 6: $\sqrt{1^2+1^2+1^2} = \sqrt{3} \approx 1.73$ 
    \item Our prediction for $K=1$ is Green, because that's the color of the closest neighbor.
    \item Our prediction for $K=3$ is Red, because it is close to two red neighbors and one green neighbors,
    so by plurality, red is the color of our prediction.
    \item If the Bayes decision boundary is highly nonlinear, then a \underline{low} K value would provide 
    the closest KNN decision boundary to Bayes. This is because low K values offer higher flexibility in the 
    KNN decision boundary, which better fits a highly nonlinear Bayes decision boundary.
\end{enumerate}
\newpage
\noindent \textbf{Question 3}
\begin{enumerate}[b.]
    \item 506 rows and 13 columns. Rows represent each data point, columns represent the entries for each data point.
    \item[c.] 
\end{enumerate}
\end{document}